% ===== 原文 pp.15–16 =====
%
% tr-02 / chunk c —— 覆盖原文 p.15–p.16（止于 §3 标题之前）；公式 tag 2.61–2.72
% ---------------------------------------------------------------------

% [原文 p.15]
\subsubsection{四点函数：显式表达式}

让我们把共形块显式确定到最低几阶。对零阶项，我们只需有
$G_{p,0,0}=l_{\alpha_1,\alpha_2,Q/2+\ii p,0}=r_{Q/2-\ii p,\alpha_3,\alpha_4,0}=1$，
于是 $F_{\alpha_1,\alpha_2,Q/2+\ii p,\alpha_3,\alpha_4}(q)=1+O(q)$。

在下一阶，我们先计算
\begin{equation}
  \langle \ee^{(Q-2\ii p)\phi(\infty)}L_1L_{-1}\ee^{(Q+2\ii p)\phi(0)}\rangle
  \tag{2.61}
\end{equation}
其中用到的是对易关系
\begin{equation}
  L_1L_{-1} = L_{-1}L_1 + 2L_0,
  \tag{2.62}
\end{equation}
再用 $L_1$ 湮灭主场 $\ee^{(Q+2\ii p)\phi}$ 这一事实。由于 $L_0$ 为 $\Delta_p=p^2+(Q/2)^2$，
我们得到
\begin{equation}
  G_{p,\{1\},\{1\}} = 2\Delta_p.
  \tag{2.63}
\end{equation}

% [Tutor Note: p.15, 式 (2.61)–(2.63)：注意 $\ee^{(Q+2\ii p)\phi}=\ee^{2\alpha\phi}$ 正是
%   $\alpha=Q/2+\ii p$ 的主场，其 $L_0=\Delta_p$；$\ee^{(Q-2\ii p)\phi}$ 是它的反射配对（$Q-\alpha$），
%   两者维数相同（见 (2.66) 的 $h_\alpha=\alpha(Q-\alpha)$ 在 $\alpha\leftrightarrow Q-\alpha$ 下不变）。
%   于是 $L_1L_{-1}|p\rangle$：$L_{-1}$ 先在 $z=0$ 造出后裔，再用 $L_1$、$L_0$ 的代数得
%   $L_1L_{-1}|p\rangle=2L_0|p\rangle=2\Delta_p|p\rangle$，即 (2.63)。这一步只用到 Virasoro 代数与
%   $L_1|p\rangle=0$，不需要知道 $C$ 的形式；(2.63) 同时是 $G^{\{1\},\{1\}}=1/(2\Delta_p)$ 的依据
%   （下面 (2.72) 会用到）。]

再下一阶，我们得到
\begin{equation}
  \begin{pmatrix}
    G_{p,\{2\},\{2\}} & G_{p,\{2\},\{1,1\}} \\
    G_{p,\{1,1\},\{2\}} & G_{p,\{1,1\},\{1,1\}}
  \end{pmatrix}
  = \begin{pmatrix}
    4\Delta_p + c/2 & 6\Delta_p \\
    6\Delta_p & 4\Delta_p + 8\Delta_0^{\,2}
  \end{pmatrix}.
  \tag{2.64}
\end{equation}

% [Tutor Note: p.15, 式 (2.64) 右下角：原文印作 $8\Delta_0^{\,2}$（已用 900 dpi 裁图复核，下标确实是 0）。
%   按 (2.56) 的 Gram 矩阵定义在能级 2 上逐项算：
%   $\langle L_2L_{-2}\rangle = \langle[L_2,L_{-2}]\rangle = 4\Delta_p+c/2$（用 (2.35)），
%   $\langle L_2L_{-1}^{2}\rangle = 3\langle L_1L_{-1}\rangle = 6\Delta_p$，
%   $\langle L_1^{2}L_{-1}^{2}\rangle = 4\Delta_p+8\Delta_p^{2}$（两次用 $L_1L_{-1}=L_{-1}L_1+2L_0$ 并注意
%   $L_1L_{-1}^{2}=L_{-1}^{2}L_1+4L_{-1}L_0+2L_{-1}$）。所以只有右下角这一处的下标与其余三处不一致，
%   应为 $\Delta_p$；正文照原文排印。]

接下来，我们需要计算
\begin{equation}
  \langle \ee^{2(Q-\alpha_1)\phi(\infty)}\ee^{2\alpha_2\phi(1)}L_{-1}\ee^{2\alpha\phi(0)}\rangle .
  \tag{2.65}
\end{equation}
为此，我们把 $\ee^{2\alpha_2\phi(1)}$ 与 $L_{-1}$ 对易。在项 $L_{-1}\ee^{2\alpha_2\phi(1)}$ 中，
$L_{-1}$ 从右侧作用到位于 $\infty$ 处的主场 $\ee^{2(Q-\alpha_1)\phi(\infty)}$ 上，把它湮灭而给出零。
对易子 $[\ee^{2\alpha_2\phi(1)},L_{-1}]$ 的计算归结为：恢复位置依赖，即把 $\phi(1)$ 换成
$\phi(z)$，对 $z$ 求导，再把 $z$ 设回 1。由共形不变性，其 $z$ 依赖恰好是
$z^{-h_{\alpha_1}+h_{\alpha_2}+h_{Q/2+\ii p}}$，其中
\begin{equation}
  h_\alpha = \alpha(Q-\alpha),\qquad
  \Delta_p = h_{Q/2+\ii p} = p^2+(Q/2)^2.
  \tag{2.66}
\end{equation}
于是我们有
\begin{equation}
  l_{\alpha_1,\alpha_2,\frac{Q}{2}+\ii p,\{1\}} = -h_{\alpha_1}+h_{\alpha_2}+\Delta_p,
  \qquad
  r_{\{1\},\frac{Q}{2}-\ii p,\alpha_3,\alpha_4} = \Delta_p+h_{\alpha_3}-h_{\alpha_4}.
  \tag{2.67}
\end{equation}

% [Tutor Note: p.15, 式 (2.67)：$l$ 的三个符号可以直接从上面那句的 $z$ 幂次读出来：
%   $z^{-h_{\alpha_1}+h_{\alpha_2}+h_{Q/2+\ii p}}$ 对 $z$ 求导等于乘上这个指数，
%   而 $\infty$、1、0 三个插入点各贡献一项，于是 $l=\big(-h_{\alpha_1}\big)+\big(h_{\alpha_2}\big)+h_{Q/2+\ii p}
%   =-h_{\alpha_1}+h_{\alpha_2}+\Delta_p$。$r$ 的符号相反（$\Delta_p+h_{\alpha_3}-h_{\alpha_4}$），
%   因为它对应另一端的主场配对；两处容易把 $-\!h_{\alpha_1}$ 与 $+\!h_{\alpha_2}$ 的符号写反。]

下一阶的项为
\begin{equation}
  l_{\alpha_1,\alpha_2,\frac{Q}{2}+\ii p,\{2\}} = -h_{\alpha_1}+2h_{\alpha_2}+\Delta_p,
  \tag{2.68}
\end{equation}
\begin{equation}
  r_{\{2\},\frac{Q}{2}-\ii p,\alpha_3,\alpha_4} = \Delta_p+2h_{\alpha_3}-h_{\alpha_4},
  \tag{2.69}
\end{equation}
\begin{equation}
  l_{\alpha_1,\alpha_2,\frac{Q}{2}+\ii p,\{1,1\}}
  = (-h_{\alpha_1}+h_{\alpha_2}+\Delta_p)(1-h_{\alpha_1}+h_{\alpha_2}+\Delta_p),
  \tag{2.70}
\end{equation}
\begin{equation}
  r_{\{1,1\},\frac{Q}{2}-\ii p,\alpha_3,\alpha_4}
  = (\Delta_p+h_{\alpha_3}-h_{\alpha_4})(1+\Delta_p+h_{\alpha_3}-h_{\alpha_4}).
  \tag{2.71}
\end{equation}

% [Tutor Note: p.15, 式 (2.68)–(2.71) 的记忆法：把 (2.67) 记作 $F=-h_{\alpha_1}+h_{\alpha_2}+\Delta_p$、
%   $G=\Delta_p+h_{\alpha_3}-h_{\alpha_4}$，则 (2.68)$=F+h_{\alpha_2}$、(2.70)$=F(F+1)$；
%   (2.69)$=G+h_{\alpha_3}$、(2.71)$=G(G+1)$。$F+1$、$G+1$ 里的"1"来自能级：第二个 $L_{-1}$
%   把有效维数抬高 1，所以两个 $L_{-1}$ 给出 $F(F+1)$ 而不是 $F^2$。抄写时最容易漏的正是这个 1。]

% [原文 p.16]
把上述结果合并，我们得到
\begin{equation}
  F_{\alpha_1,\alpha_2,Q/2+\ii p,\alpha_3,\alpha_4}(q)
  = 1 + \frac{(-h_{\alpha_1}+h_{\alpha_2}+\Delta_p)(\Delta_p+h_{\alpha_3}-h_{\alpha_4})}{2\Delta_p}
    \,q + O(q^2).
  \tag{2.72}
\end{equation}
$q^2$ 阶的项可以由上面给出的数据算出，但太长，不便在此列出。

% [Tutor Note: p.16, 式 (2.72)：用 (2.60) 验证一遍——一阶项只需 $\{n\}=\{n'\}=\{1\}$，
%   $F = q\cdot l_{\cdots,\{1\}}\cdot G^{\{1\},\{1\}}\cdot r_{\{1\},\cdots}$，而由 (2.63)
%   $G^{\{1\},\{1\}}=1/G_{p,\{1\},\{1\}}=1/(2\Delta_p)$，代入 (2.67) 即得
%   $q\cdot(-h_{\alpha_1}+h_{\alpha_2}+\Delta_p)(\Delta_p+h_{\alpha_3}-h_{\alpha_4})/(2\Delta_p)$，
%   与 (2.72) 完全一致。这条自检同时说明了 (2.60) 里 $G$ 用逆矩阵、$q$ 的幂次取 $\sum n$ 这两处约定。]

要得到共形块 $F(q)$ 中 $q$ 的更高阶项，工作是繁琐的，但并不困难。注意，这完全由 Virasoro
代数的性质决定，而最终结果只用 $h_{\alpha_i}=\alpha_i(Q-\alpha_i)$ 与 $\Delta_p$
（它们分别是主场 $\ee^{2\alpha_i\phi}$ 与 $\ee^{(Q+2\ii p)\phi}$ 的 $L_0$）以及中心荷 $c$ 表示。
此时，在读者自选的计算机代数系统里写一个程序、把共形块算到 $q$ 的任意阶，是很有启发性的。

% [Tutor Note: p.16, "完全由 Virasoro 代数的性质决定"：这句话是 §2.2 的技术总结——(2.72) 型的结果
%   只依赖中心荷 $c$、各主场最高权 $h_{\alpha_i}$ 与中间态最高权 $\Delta_p$ 这三类数据，
%   与 Liouville 的拉格朗日量、$\mu$、$b$ 的具体取值无关。这正是后文 §5 能用 Liouville 一侧去
%   "匹配" 四维规范理论一侧配分函数的技术前提：两侧的共形块是同一条 Virasoro 递推。]
